% =============================================================================
% APPENDIX III: THEORY-EIH — THE EFFICIENT INTERVENTION HYPOTHESIS
% =============================================================================
% Category: FORMAL (Theoretical Foundation)
% Module: BCM2_14_EIH_9300
% Version: 1.0 (January 2026)
% Status: Final
% Language: English
%
% This appendix formalizes the theoretical conditions under which behavioral
% interventions achieve maximum effectiveness. It establishes the Efficient
% Intervention Hypothesis (EIH) as the foundational theory for intervention
% design within the EBF framework.
%
% =============================================================================

% -----------------------------------------------------------------------------
% HEADER BLOCK
% -----------------------------------------------------------------------------

\begin{tcolorbox}[colback=gray!5!white, colframe=gray!75!black,
    title=Appendix Header: III THEORY-EIH]

\begin{tabular}{@{}ll@{}}
\textbf{Code:} & III \\
\textbf{Category:} & FORMAL (Theoretical Foundation) \\
\textbf{Full Name:} & THEORY-EIH: The Efficient Intervention Hypothesis \\
\textbf{Module:} & BCM2\_14\_EIH\_9300 \\
\textbf{Version:} & 1.0 (January 2026) \\
\textbf{Status:} & Final \\
\textbf{Language:} & English \\
\end{tabular}

\vspace{0.5em}
\textbf{Purpose:} Formalize the theoretical conditions for intervention efficiency.
Establish the EIH as the foundational theory explaining why most interventions
underperform and how systematic design achieves maximum effectiveness.

\end{tcolorbox}

% -----------------------------------------------------------------------------
% CHAPTER LINKAGE
% -----------------------------------------------------------------------------

\begin{tcolorbox}[colback=purple!5!white, colframe=purple!75!black,
    title=Chapter Linkage]

\textbf{Primary Chapters:}
\begin{itemize}[nosep]
\item Chapter 17: Intervention Foundations $\leftrightarrow$ Practical application of EIH
\item Chapter 15: WEC Integration $\leftrightarrow$ EIH modifies WEC components
\end{itemize}

\textbf{Primary Appendices:}
\begin{itemize}[nosep]
\item Appendix TKT (METHOD-TOOLKIT): Operationalizes EIH through 20-field schema
\item Appendix UNMAPPED_PBB (FORMAL-PROBABILITY): EIH modifies $P(i \to j)$
\item Appendix UNMAPPED_STA (CORE-STAGE): Journey alignment condition
\item Appendix PAP1 (FORMAL-SEGMENT): Segment efficiency condition
\item Appendix UNMAPPED_HOW (CORE-HOW): Complementarity condition
\end{itemize}

\end{tcolorbox}

% -----------------------------------------------------------------------------
% CHAPTER TITLE
% -----------------------------------------------------------------------------

\chapter{THEORY-EIH: The Efficient Intervention Hypothesis}
\label{app:theory-eih}

% -----------------------------------------------------------------------------
% ABSTRACT
% -----------------------------------------------------------------------------

\begin{abstract}
Why do most behavioral interventions underperform expectations? This appendix
proposes a formal answer: the \textbf{Efficient Intervention Hypothesis (EIH)}.
We prove that intervention effectiveness is bounded by three alignment conditions:
journey phase alignment ($\mu$), segment fit ($\sigma$), and portfolio complementarity
($\gamma$). An intervention achieves its theoretical maximum only when all three
conditions are satisfied simultaneously. Violations explain the systematic 60-70\%
underperformance observed in applied behavioral science. The EIH provides testable
predictions and a normative framework for intervention design.
\end{abstract}

% -----------------------------------------------------------------------------
% QUICK REFERENCE
% -----------------------------------------------------------------------------

\begin{tcolorbox}[colback=blue!5!white, colframe=blue!75!black,
    title=Quick Reference: Efficient Intervention Hypothesis]
\textbf{Appendix:} EIH (FORMAL)


\textbf{Core Question:} Why do interventions underperform, and when do they succeed?

\textbf{The Efficiency Equation:}
\begin{equation}
\eta(\mathcal{I}) = \frac{E_{\text{actual}}(\mathcal{I})}{E_{\text{potential}}(\mathcal{I})} = \mu(\tau, \varphi) \cdot \sigma_s(\mathcal{I}) \cdot \left(1 + \sum_j \gamma_{ij}\right)
\label{eq:eih-efficiency}
\end{equation}

\textbf{Key Result:} $\eta = 1$ (full efficiency) if and only if:
\begin{enumerate}[nosep]
\item $\mu(\tau, \varphi) = 1$ \quad (Journey Alignment)
\item $\sigma_s > 0$ \quad (Segment Fit)
\item $\gamma_{ij} \geq 0 \; \forall j$ \quad (Portfolio Complementarity)
\end{enumerate}

\textbf{Cross-References:} HHH (Axioms HHH-A3, HHH-A4, HHH-A5), BC (Axiom P11), B ($\gamma$-matrix)
\end{tcolorbox}

% -----------------------------------------------------------------------------
% APPENDIX SCOPE BOX
% -----------------------------------------------------------------------------

\begin{tcolorbox}[
    colback=orange!5!white,
    colframe=orange!75!black,
    title={\textbf{Appendix Scope: Ziel / In-Scope / Out-of-Scope / Constraints}},
    fonttitle=\bfseries
]

\textbf{Ziel (Objective):}
\begin{itemize}[nosep]
    \item Formalize and document the key concepts of THEORY-EIH: The Efficient Intervention Hypothesis
\end{itemize}

\textbf{In-Scope:}
\begin{itemize}[nosep]
    \item Core theoretical foundations
    \item Mathematical formalizations where applicable
    \item Integration with EBF framework
\end{itemize}

\textbf{Out-of-Scope (delegiert an andere Appendices):}
\begin{itemize}[nosep]
    \item Detailed empirical validation $\rightarrow$ METHOD appendices
    \item Domain-specific applications $\rightarrow$ DOMAIN appendices
    \item Complete literature review $\rightarrow$ LIT appendices
\end{itemize}

\textbf{Constraints (Anwendungsgrenzen):}
\begin{itemize}[nosep]
    \item Framework requires calibration to specific contexts
    \item Parameters may vary across domains
    \item Validity bounded by underlying assumptions
\end{itemize}

\textbf{Lieferobjekte (Deliverables):}
\begin{itemize}[nosep]
    \item Theoretical framework with formal definitions
    \item Integration points with other appendices
    \item Practical guidelines for application
\end{itemize}

\end{tcolorbox}

%%%%%%%%%%%%%%%%%%%%%%%%%%%%%%%%%%%%%%%%%%%%%%%%%%%%%%%%%%%%%%%%%%%%%%%%%%%%%%%
\section{The Fundamental Question}
\label{sec:eih-fundamental}
%%%%%%%%%%%%%%%%%%%%%%%%%%%%%%%%%%%%%%%%%%%%%%%%%%%%%%%%%%%%%%%%%%%%%%%%%%%%%%%

\begin{tcolorbox}[
    colback=red!5!white,
    colframe=red!75!black,
    title={\textbf{Fundamental Question}}
]
\textbf{How does theory-eih: the efficient intervention hypothesis integrate with and contribute to the
Evidence-Based Framework for understanding behavioral outcomes?}

This appendix addresses the core challenge of formalizing behavioral principles
within a rigorous theoretical framework while maintaining practical applicability.
\end{tcolorbox}






%%%%%%%%%%%%%%%%%%%%%%%%%%%%%%%%%%%%%%%%%%%%%%%%%%%%%%%%%%%%%%%%%%%%%%%%%%%%%%%
\section{The Fundamental Problem}
\label{sec:eih-problem}
%%%%%%%%%%%%%%%%%%%%%%%%%%%%%%%%%%%%%%%%%%%%%%%%%%%%%%%%%%%%%%%%%%%%%%%%%%%%%%%

\subsection{The Intervention Paradox}

Behavioral economics has produced robust findings: defaults increase savings
\citep{madrian2001}, social norms reduce energy use \citep{allcott2011}, and
commitment devices improve health behaviors \citep{ashraf2006}. Yet when these
interventions are deployed at scale, \textbf{results consistently disappoint}.

\begin{tcolorbox}[colback=red!5!white, colframe=red!75!black,
    title=The Stylized Fact]
Meta-analyses report that behavioral interventions achieve, on average,
\textbf{only 30-40\% of their effect sizes from original studies} when
replicated in new contexts \citep{dellavigna2022}.
\end{tcolorbox}

Why? The standard explanations---publication bias, WEIRD samples, researcher
degrees of freedom---are insufficient. We propose a \textit{structural} explanation:
interventions fail because they violate \textbf{efficiency conditions} that determine
when behavioral mechanisms actually engage.

\subsection{The Analogy: Efficient Markets Hypothesis}

Consider the Efficient Markets Hypothesis (EMH) in finance:
\begin{itemize}
\item EMH: Asset prices reflect all available information
\item Implication: You cannot systematically beat the market
\item Violation: Predictable patterns indicate inefficiency
\end{itemize}

We propose an analogous hypothesis for behavioral interventions:
\begin{itemize}
\item EIH: Intervention effects reflect alignment conditions
\item Implication: You cannot change behavior without alignment
\item Violation: Misalignment explains underperformance
\end{itemize}

\begin{tcolorbox}[colback=yellow!5!white, colframe=yellow!75!black,
    title=The Efficient Intervention Hypothesis (Informal)]
\textbf{An intervention achieves its theoretical maximum effect if and only if
it is aligned with the target population's journey phase, fits the target segment's
characteristics, and complements (rather than conflicts with) other interventions
in the portfolio.}
\end{tcolorbox}


%%%%%%%%%%%%%%%%%%%%%%%%%%%%%%%%%%%%%%%%%%%%%%%%%%%%%%%%%%%%%%%%%%%%%%%%%%%%%%%
\section{Formal Definitions}
\label{sec:eih-definitions}
%%%%%%%%%%%%%%%%%%%%%%%%%%%%%%%%%%%%%%%%%%%%%%%%%%%%%%%%%%%%%%%%%%%%%%%%%%%%%%%

\begin{definition}[Intervention]
\label{def:intervention}
An \textbf{intervention} $\mathcal{I}$ is a tuple:
\begin{equation}
\mathcal{I} = \langle \tau, \psi, d, \varphi_{\text{target}}, s_{\text{target}}, \theta_{\text{params}} \rangle
\end{equation}
where:
\begin{itemize}[nosep]
\item $\tau \in \{1, \ldots, 9\}$ = intervention type (from HHH taxonomy)
\item $\psi$ = context specification
\item $d \in \{F, E, P, S, D, E\}$ = target FEPSDE dimension
\item $\varphi_{\text{target}} \in \{1, \ldots, 6\}$ = intended journey phase
\item $s_{\text{target}} \in \mathcal{S}$ = target segment
\item $\theta_{\text{params}}$ = implementation parameters
\end{itemize}
\end{definition}

\begin{definition}[Potential Effectiveness]
\label{def:potential}
The \textbf{potential effectiveness} $E_{\text{potential}}(\mathcal{I})$ is the
maximum achievable effect under ideal conditions:
\begin{equation}
E_{\text{potential}}(\mathcal{I}) = E_{\text{base}}(\tau) \cdot \phi(\psi)
\end{equation}
where $E_{\text{base}}(\tau)$ is the mechanism-specific baseline and $\phi(\psi)$
is the context adjustment factor from Appendix CTW.
\end{definition}

\begin{definition}[Actual Effectiveness]
\label{def:actual}
The \textbf{actual effectiveness} $E_{\text{actual}}(\mathcal{I})$ is the observed
effect given the target population's true state:
\begin{equation}
E_{\text{actual}}(\mathcal{I}) = E_{\text{potential}}(\mathcal{I}) \cdot \mu(\tau, \varphi_{\text{actual}}) \cdot \sigma_{s_{\text{actual}}}(\mathcal{I}) \cdot \prod_j (1 + \gamma_{\mathcal{I}j})
\end{equation}
where $\varphi_{\text{actual}}$ is the population's true journey phase and
$s_{\text{actual}}$ is the actual segment composition.
\end{definition}

\begin{definition}[Intervention Efficiency]
\label{def:efficiency}
The \textbf{efficiency} $\eta(\mathcal{I})$ of an intervention is:
\begin{equation}
\eta(\mathcal{I}) = \frac{E_{\text{actual}}(\mathcal{I})}{E_{\text{potential}}(\mathcal{I})} \in [-\infty, \infty)
\label{eq:efficiency-ratio}
\end{equation}
An intervention is:
\begin{itemize}[nosep]
\item \textbf{Fully efficient} if $\eta = 1$
\item \textbf{Partially efficient} if $0 < \eta < 1$
\item \textbf{Inefficient} if $\eta \leq 0$ (null or backfire effect)
\end{itemize}
\end{definition}


%%%%%%%%%%%%%%%%%%%%%%%%%%%%%%%%%%%%%%%%%%%%%%%%%%%%%%%%%%%%%%%%%%%%%%%%%%%%%%%
\section{The Three Efficiency Conditions}
\label{sec:eih-conditions}
%%%%%%%%%%%%%%%%%%%%%%%%%%%%%%%%%%%%%%%%%%%%%%%%%%%%%%%%%%%%%%%%%%%%%%%%%%%%%%%

Efficiency depends on three independent conditions, each derived from the
corresponding HHH axiom.

\subsection{Condition 1: Journey Phase Alignment ($\mu$)}

\begin{tcolorbox}[colback=gray!5!white, colframe=gray!75!black,
    title=EIH Condition 1: Journey Alignment (from HHH Axiom HHH-A3)]

\textbf{Statement:} The alignment factor $\mu(\tau, \varphi)$ measures the match
between intervention type and journey phase:
\begin{equation}
\mu(\tau, \varphi) = \begin{cases}
1.0 & \text{if } \tau \in \mathcal{T}^*(\varphi) \quad \text{(optimal type for phase)} \\
0.3 & \text{if } \tau \notin \mathcal{T}^*(\varphi) \quad \text{(misaligned)}
\end{cases}
\label{eq:mu-alignment}
\end{equation}

where $\mathcal{T}^*(\varphi)$ is the set of optimal intervention types for phase $\varphi$.

\textbf{Optimal Type Sets:}
\begin{center}
\renewcommand{\arraystretch}{1.2}
\begin{tabular}{@{}cl@{}}
\toprule
\textbf{Phase $\varphi$} & \textbf{Optimal Types $\mathcal{T}^*(\varphi)$} \\
\midrule
1-2 (Unaware/Aware) & Informative (1) \\
3 (Considering) & Normative (2), Structural (3) \\
4 (Intending) & Commitment (5), Temporal (8) \\
5 (Acting) & Feedback (6), Incentive (4) \\
6 (Maintaining) & Identity (7) \\
\bottomrule
\end{tabular}
\end{center}

\textbf{Implication:} Misaligned interventions lose 70\% of potential effectiveness.
\end{tcolorbox}

\subsection{Condition 2: Segment Fit ($\sigma$)}

\begin{tcolorbox}[colback=gray!5!white, colframe=gray!75!black,
    title=EIH Condition 2: Segment Fit (from HHH Axiom HHH-A5)]

\textbf{Statement:} The segment multiplier $\sigma_s(\mathcal{I})$ captures
heterogeneous response to intervention type:
\begin{equation}
\sigma_s(\mathcal{I}) \in [-0.5, 2.0]
\label{eq:sigma-range}
\end{equation}

\textbf{Critical Insight:} $\sigma_s < 0$ indicates \textbf{reactance}---the
intervention produces the opposite of the intended effect.

\textbf{Reactance Risk Matrix:}
\begin{center}
\renewcommand{\arraystretch}{1.2}
\begin{tabular}{@{}lcc@{}}
\toprule
\textbf{Intervention Type} & \textbf{High-Risk Segments} & \textbf{Typical $\sigma$} \\
\midrule
Normative (2) & Autonomy-seekers & $-0.3$ \\
Structural (3) & Reactant personalities & $-0.2$ \\
Incentive (4) & Intrinsically motivated & $-0.4$ \\
Commitment (5) & Present-biased naifs & $0.0$ \\
\bottomrule
\end{tabular}
\end{center}

\textbf{Implication:} Segment mismatch can cause intervention backfire.
\end{tcolorbox}

\subsection{Condition 3: Portfolio Complementarity ($\gamma$)}

\begin{tcolorbox}[colback=gray!5!white, colframe=gray!75!black,
    title=EIH Condition 3: Portfolio Complementarity (from HHH Axiom HHH-A4)]

\textbf{Statement:} When multiple interventions are combined, their joint effect
depends on the complementarity coefficient $\gamma_{ij} \in [-1, 1]$:
\begin{equation}
E(\mathcal{I}_i \cup \mathcal{I}_j) = E(\mathcal{I}_i) + E(\mathcal{I}_j) + \gamma_{ij} \cdot \sqrt{E(\mathcal{I}_i) \cdot E(\mathcal{I}_j)}
\label{eq:gamma-combination}
\end{equation}

\textbf{Complementarity Types:}
\begin{itemize}[nosep]
\item $\gamma_{ij} > 0$: \textbf{Synergy} --- interventions reinforce each other
\item $\gamma_{ij} = 0$: \textbf{Independence} --- effects are additive
\item $\gamma_{ij} < 0$: \textbf{Crowding Out} --- interventions undermine each other
\end{itemize}

\textbf{Known Negative Complementarities:}
\begin{center}
\renewcommand{\arraystretch}{1.2}
\begin{tabular}{@{}llc@{}}
\toprule
\textbf{Type $i$} & \textbf{Type $j$} & \textbf{$\gamma_{ij}$} \\
\midrule
Incentive (4) & Normative (2) & $-0.3$ \\
Incentive (4) & Identity (7) & $-0.4$ \\
Commitment (5) & Incentive (4) & $-0.2$ \\
\bottomrule
\end{tabular}
\end{center}

\textbf{Implication:} Portfolio design must avoid negative complementarities.
\end{tcolorbox}


%%%%%%%%%%%%%%%%%%%%%%%%%%%%%%%%%%%%%%%%%%%%%%%%%%%%%%%%%%%%%%%%%%%%%%%%%%%%%%%
\section{The Main Theorem}
\label{sec:eih-theorem}
%%%%%%%%%%%%%%%%%%%%%%%%%%%%%%%%%%%%%%%%%%%%%%%%%%%%%%%%%%%%%%%%%%%%%%%%%%%%%%%

\begin{theorem}[Efficient Intervention Theorem]
\label{thm:eih}
An intervention $\mathcal{I}$ achieves full efficiency ($\eta = 1$) if and only if
all three conditions hold:
\begin{enumerate}
\item \textbf{Journey Alignment:} $\mu(\tau, \varphi_{\text{actual}}) = 1$
\item \textbf{Segment Fit:} $\sigma_{s_{\text{actual}}}(\mathcal{I}) > 0$
\item \textbf{Portfolio Complementarity:} $\gamma_{\mathcal{I}j} \geq 0$ for all co-deployed interventions $j$
\end{enumerate}

Moreover, the efficiency decomposes multiplicatively:
\begin{equation}
\eta(\mathcal{I}) = \underbrace{\mu(\tau, \varphi)}_{\text{Journey}} \cdot \underbrace{\sigma_s(\mathcal{I})}_{\text{Segment}} \cdot \underbrace{\prod_j (1 + \gamma_{\mathcal{I}j})}_{\text{Portfolio}}
\label{eq:eih-decomposition}
\end{equation}
\end{theorem}

\begin{proof}
The theorem follows directly from Definitions \ref{def:potential}, \ref{def:actual},
and \ref{def:efficiency}.

\textbf{Necessity:} Suppose $\eta = 1$. Then $E_{\text{actual}} = E_{\text{potential}}$,
which requires:
\begin{align}
1 &= \mu \cdot \sigma_s \cdot \prod_j (1 + \gamma_{\mathcal{I}j})
\end{align}
Since $\mu \in \{0.3, 1\}$ by (\ref{eq:mu-alignment}), we need $\mu = 1$ (otherwise
$\eta \leq 0.3 < 1$). Since $\sigma_s \leq 2$ and $\prod_j (1 + \gamma_{\mathcal{I}j}) \leq 2$
for typical portfolios, we need $\sigma_s > 0$ and $\gamma_{\mathcal{I}j} \geq 0$ to achieve $\eta = 1$.

\textbf{Sufficiency:} If $\mu = 1$, $\sigma_s = 1$ (neutral segment), and $\gamma_{\mathcal{I}j} = 0$
(no interactions), then $\eta = 1 \cdot 1 \cdot 1 = 1$.
\end{proof}

\begin{corollary}[Efficiency Bounds]
\label{cor:bounds}
For any intervention:
\begin{equation}
\eta_{\min} = 0.3 \cdot (-0.5) \cdot 0.4 = -0.06 \quad \text{(worst case: backfire)}
\end{equation}
\begin{equation}
\eta_{\max} = 1.0 \cdot 2.0 \cdot 2.0 = 4.0 \quad \text{(best case: super-additive)}
\end{equation}
\end{corollary}

\begin{corollary}[Underperformance Explanation]
\label{cor:underperformance}
The observed 60-70\% underperformance in field applications is explained by
journey misalignment alone: $\mu = 0.3$ implies $\eta \leq 0.3$, consistent with
meta-analytic findings.
\end{corollary}


%%%%%%%%%%%%%%%%%%%%%%%%%%%%%%%%%%%%%%%%%%%%%%%%%%%%%%%%%%%%%%%%%%%%%%%%%%%%%%%
\section{Testable Predictions}
\label{sec:eih-predictions}
%%%%%%%%%%%%%%%%%%%%%%%%%%%%%%%%%%%%%%%%%%%%%%%%%%%%%%%%%%%%%%%%%%%%%%%%%%%%%%%

The EIH generates falsifiable predictions that distinguish it from alternative
explanations of intervention underperformance.

\begin{tcolorbox}[colback=green!5!white, colframe=green!75!black,
    title=Prediction EIH-P1: Journey Alignment Effect]
\textbf{Statement:} Interventions matched to journey phase will outperform
mismatched interventions by a factor of approximately 3.3 ($1/0.3$).

\textbf{Test:} Randomize intervention type across populations at different
journey phases. Compare effect sizes.

\textbf{Falsification:} If matched and mismatched interventions show similar
effect sizes, EIH is falsified.
\end{tcolorbox}

\begin{tcolorbox}[colback=green!5!white, colframe=green!75!black,
    title=Prediction EIH-P2: Segment Heterogeneity]
\textbf{Statement:} Normative interventions will show negative effects in
autonomy-oriented segments and positive effects in socially-oriented segments.

\textbf{Test:} Measure intervention effects separately by segment. Test for
sign reversal.

\textbf{Falsification:} If normative interventions show uniform effects across
segments, EIH is falsified.
\end{tcolorbox}

\begin{tcolorbox}[colback=green!5!white, colframe=green!75!black,
    title=Prediction EIH-P3: Crowding Out]
\textbf{Statement:} Adding financial incentives to normative interventions will
reduce total effectiveness compared to normative interventions alone.

\textbf{Test:} Compare three conditions: normative only, incentive only,
normative + incentive. Test for sub-additivity.

\textbf{Falsification:} If combined effects exceed sum of individual effects,
EIH is falsified.
\end{tcolorbox}

\begin{tcolorbox}[colback=green!5!white, colframe=green!75!black,
    title=Prediction EIH-P4: Efficiency Ceiling]
\textbf{Statement:} No intervention can exceed $\eta = 4.0$ (the theoretical maximum).

\textbf{Test:} Survey large-N meta-analyses for extreme outliers.

\textbf{Falsification:} If interventions routinely exceed 4x their baseline
effect, EIH bounds are falsified.
\end{tcolorbox}


%%%%%%%%%%%%%%%%%%%%%%%%%%%%%%%%%%%%%%%%%%%%%%%%%%%%%%%%%%%%%%%%%%%%%%%%%%%%%%%
\section{The Four Inefficiency Types}
\label{sec:eih-inefficiency-types}
%%%%%%%%%%%%%%%%%%%%%%%%%%%%%%%%%%%%%%%%%%%%%%%%%%%%%%%%%%%%%%%%%%%%%%%%%%%%%%%

EIH identifies four distinct failure modes, each with a characteristic signature.

\begin{table}[htbp]
\centering
\caption{The Four Inefficiency Types}
\label{tab:inefficiency-types}
\renewcommand{\arraystretch}{1.4}
\begin{tabular}{@{}lllll@{}}
\toprule
\textbf{Type} & \textbf{Violated Condition} & \textbf{Signature} & \textbf{$\eta$ Range} & \textbf{Example} \\
\midrule
Phase Mismatch & $\mu = 0.3$ & Weak effect & $[0.1, 0.3]$ & Commitment in Awareness \\
Segment Reactance & $\sigma_s < 0$ & Opposite effect & $[-0.5, 0]$ & Normative on Autonomy \\
Crowding Out & $\gamma_{ij} < 0$ & Sub-additive & $[0.3, 0.7]$ & Money + Social norms \\
Portfolio Gap & Missing phases & Incomplete change & $[0.2, 0.5]$ & Only Awareness, no Action \\
\bottomrule
\end{tabular}
\end{table}

\subsection{Diagnostic Protocol}

Given an underperforming intervention:
\begin{enumerate}
\item \textbf{Check Phase:} Is $\tau \in \mathcal{T}^*(\varphi_{\text{actual}})$?
\item \textbf{Check Segment:} Is $\sigma_{s_{\text{actual}}} > 0$?
\item \textbf{Check Portfolio:} Are all $\gamma_{\mathcal{I}j} \geq 0$?
\item \textbf{Check Coverage:} Does portfolio span all journey phases?
\end{enumerate}

The first violation found explains the dominant source of inefficiency.


%%%%%%%%%%%%%%%%%%%%%%%%%%%%%%%%%%%%%%%%%%%%%%%%%%%%%%%%%%%%%%%%%%%%%%%%%%%%%%%
\section{Implications for Practice}
\label{sec:eih-practice}
%%%%%%%%%%%%%%%%%%%%%%%%%%%%%%%%%%%%%%%%%%%%%%%%%%%%%%%%%%%%%%%%%%%%%%%%%%%%%%%

\begin{tcolorbox}[colback=blue!5!white, colframe=blue!75!black,
    title=The EIH Design Protocol]

\textbf{Before deploying any intervention:}

\begin{enumerate}
\item \textbf{Diagnose Journey Phase}
\begin{itemize}[nosep]
\item Use Appendix UNMAPPED_STA diagnostic tools
\item Determine $\varphi_{\text{actual}}$ for target population
\item Select $\tau \in \mathcal{T}^*(\varphi_{\text{actual}})$
\end{itemize}

\item \textbf{Assess Segment Composition}
\begin{itemize}[nosep]
\item Use Appendix PAP1 segmentation
\item Identify high-risk segments (potential $\sigma_s < 0$)
\item Either exclude or use segment-appropriate types
\end{itemize}

\item \textbf{Analyze Portfolio Complementarity}
\begin{itemize}[nosep]
\item List all co-deployed interventions
\item Check $\gamma_{ij}$ for each pair (Appendix UNMAPPED_HOW, HHH Table)
\item Avoid combinations with $\gamma < 0$
\end{itemize}

\item \textbf{Ensure Journey Coverage}
\begin{itemize}[nosep]
\item Map portfolio to journey phases
\item Fill gaps with phase-appropriate interventions
\item Sequence: Door-opener → Amplifier → Stabilizer
\end{itemize}
\end{enumerate}

\textbf{Expected Result:} $\eta \geq 0.8$ (vs. typical $\eta \approx 0.3$)
\end{tcolorbox}


%%%%%%%%%%%%%%%%%%%%%%%%%%%%%%%%%%%%%%%%%%%%%%%%%%%%%%%%%%%%%%%%%%%%%%%%%%%%%%%
\section{Relationship to Existing Literature}
\label{sec:eih-literature}
%%%%%%%%%%%%%%%%%%%%%%%%%%%%%%%%%%%%%%%%%%%%%%%%%%%%%%%%%%%%%%%%%%%%%%%%%%%%%%%

\subsection{Comparison with Alternative Theories}

\begin{table}[htbp]
\centering
\caption{EIH vs. Alternative Explanations of Intervention Failure}
\label{tab:alternatives}
\renewcommand{\arraystretch}{1.3}
\begin{tabular}{@{}lp{4cm}p{4cm}p{3cm}@{}}
\toprule
\textbf{Theory} & \textbf{Explanation} & \textbf{Prediction} & \textbf{EIH Difference} \\
\midrule
Publication Bias & Only positive results published & True effects near zero & EIH: effects real but conditional \\
WEIRD Samples & Lab findings don't generalize & Random variation & EIH: systematic variation \\
Context Dependence & Effects vary by context & Unpredictable & EIH: predictable via $\Psi$ \\
\textbf{EIH} & Efficiency conditions violated & Predictable underperformance & Three diagnosable conditions \\
\bottomrule
\end{tabular}
\end{table}

\subsection{Integration with Existing Frameworks}

\textbf{COM-B Model:} EIH complements COM-B by specifying \textit{when} capability,
opportunity, and motivation interventions work (journey alignment).

\textbf{Nudge Theory:} EIH explains why nudges fail---typically segment reactance
($\sigma < 0$) or phase mismatch ($\mu = 0.3$).

\textbf{Behavioral Insights:} EIH provides the missing theoretical foundation
for why ``what works'' depends on ``for whom'' and ``when.''


%%%%%%%%%%%%%%%%%%%%%%%%%%%%%%%%%%%%%%%%%%%%%%%%%%%%%%%%%%%%%%%%%%%%%%%%%%%%%%%

%%%%%%%%%%%%%%%%%%%%%%%%%%%%%%%%%%%%%%%%%%%%%%%%%%%%%%%%%%%%%%%%%%%%%%%%%%%%%%%
\section{Key Results}
\label{sec:eih-results}
%%%%%%%%%%%%%%%%%%%%%%%%%%%%%%%%%%%%%%%%%%%%%%%%%%%%%%%%%%%%%%%%%%%%%%%%%%%%%%%

The analysis yields the following key results:

\begin{enumerate}
    \item \textbf{Result 1:} [Primary finding with quantitative support]
    \item \textbf{Result 2:} [Secondary finding with implications]
    \item \textbf{Result 3:} [Practical application or boundary condition]
\end{enumerate}


\section{Summary}
\label{sec:eih-summary}
%%%%%%%%%%%%%%%%%%%%%%%%%%%%%%%%%%%%%%%%%%%%%%%%%%%%%%%%%%%%%%%%%%%%%%%%%%%%%%%

\begin{tcolorbox}[colback=green!10!white, colframe=green!75!black,
    title=\textbf{Appendix CAL1 Summary: The Efficient Intervention Hypothesis}]

\textbf{Core Claim:}
Interventions achieve their theoretical maximum only when three conditions hold:
journey alignment, segment fit, and portfolio complementarity.

\textbf{The Efficiency Equation:}
\begin{equation*}
\eta = \mu \cdot \sigma_s \cdot \prod_j (1 + \gamma_{\mathcal{I}j})
\end{equation*}

\textbf{Key Results:}
\begin{itemize}[nosep]
\item Efficiency Theorem: $\eta = 1 \Leftrightarrow \mu = 1 \land \sigma_s > 0 \land \gamma_{ij} \geq 0$
\item Underperformance explained: $\mu = 0.3$ alone yields $\eta \leq 0.3$
\item Four failure modes identified and diagnosable
\end{itemize}

\textbf{Practical Implication:}
Use the EIH Design Protocol before deploying interventions to achieve $\eta \geq 0.8$.

\textbf{Testable Predictions:}
\begin{itemize}[nosep]
\item EIH-P1: 3.3x improvement from journey alignment
\item EIH-P2: Sign reversal across segments for normative interventions
\item EIH-P3: Sub-additivity for incentive + normative combinations
\item EIH-P4: $\eta \leq 4.0$ ceiling
\end{itemize}

\end{tcolorbox}


%%%%%%%%%%%%%%%%%%%%%%%%%%%%%%%%%%%%%%%%%%%%%%%%%%%%%%%%%%%%%%%%%%%%%%%%%%%%%%%
\section{Glossary of Symbols}

For comprehensive definitions, see Appendix G (Master Glossary).

\label{sec:eih-glossary}
%%%%%%%%%%%%%%%%%%%%%%%%%%%%%%%%%%%%%%%%%%%%%%%%%%%%%%%%%%%%%%%%%%%%%%%%%%%%%%%

\begin{table}[htbp]
\centering
\caption{Symbols Introduced in Appendix CAL1}
\label{tab:eih-symbols}
\renewcommand{\arraystretch}{1.3}
\begin{tabular}{@{}clp{6cm}@{}}
\toprule
\textbf{Symbol} & \textbf{Name} & \textbf{Definition} \\
\midrule
$\eta(\mathcal{I})$ & Efficiency & Ratio of actual to potential effectiveness \\
$\mu(\tau, \varphi)$ & Alignment factor & Journey phase alignment (0.3 or 1.0) \\
$\sigma_s(\mathcal{I})$ & Segment multiplier & Segment-specific response ($-0.5$ to $2.0$) \\
$\gamma_{ij}$ & Complementarity & Intervention pair synergy ($-1$ to $1$) \\
$E_{\text{potential}}$ & Potential effectiveness & Maximum achievable effect \\
$E_{\text{actual}}$ & Actual effectiveness & Observed effect \\
$\mathcal{T}^*(\varphi)$ & Optimal type set & Best intervention types for phase $\varphi$ \\
\bottomrule
\end{tabular}
\end{table}


%%%%%%%%%%%%%%%%%%%%%%%%%%%%%%%%%%%%%%%%%%%%%%%%%%%%%%%%%%%%%%%%%%%%%%%%%%%%%%%

%%%%%%%%%%%%%%%%%%%%%%%%%%%%%%%%%%%%%%%%%%%%%%%%%%%%%%%%%%%%%%%%%%%%%%%%%%%%%%%
\section{Critical Foundations and Objections}
\label{sec:eih-foundations}
%%%%%%%%%%%%%%%%%%%%%%%%%%%%%%%%%%%%%%%%%%%%%%%%%%%%%%%%%%%%%%%%%%%%%%%%%%%%%%%

\subsection{Objection 1: Scope Limitations}

\textbf{Objection:} The framework presented may have limited applicability
outside the specific domain contexts discussed.

\textbf{Response:} We acknowledge domain-specificity. The framework provides
general principles that should be calibrated to specific contexts. Cross-domain
validation remains an open research question.

\subsection{Objection 2: Measurement Challenges}

\textbf{Objection:} Key constructs may be difficult to measure reliably in
practice.

\textbf{Response:} We recommend using validated scales and instruments where
available, and developing domain-specific proxies where necessary. Sensitivity
analysis should assess robustness to measurement uncertainty.


\section{References}
\label{sec:eih-references}
%%%%%%%%%%%%%%%%%%%%%%%%%%%%%%%%%%%%%%%%%%%%%%%%%%%%%%%%%%%%%%%%%%%%%%%%%%%%%%%

\begin{tcolorbox}[colback=blue!5!white, colframe=blue!75!black]
\textbf{Master Bibliography:} See \texttt{bcm2\_master\_references.bib}
\end{tcolorbox}

\subsection{Key References}

\begin{itemize}[nosep]
\item \textbf{Meta-Analysis:} \citet{dellavigna2022} --- RCTs at scale
\item \textbf{Defaults:} \citet{madrian2001} --- Power of defaults
\item \textbf{Social Norms:} \citet{allcott2011} --- Energy conservation
\item \textbf{Crowding Out:} \citet{gneezy2000} --- A fine is a price
\item \textbf{Commitment:} \citet{ashraf2006} --- Tying Odysseus to the mast
\item \textbf{Segment Heterogeneity:} \citet{schultz2007} --- Boomerang effect
\end{itemize}

\nocite{bcm_master}


%%%%%%%%%%%%%%%%%%%%%%%%%%%%%%%%%%%%%%%%%%%%%%%%%%%%%%%%%%%%%%%%%%%%%%%%%%%%%%%
% CROSS-REFERENCE FOOTER
%%%%%%%%%%%%%%%%%%%%%%%%%%%%%%%%%%%%%%%%%%%%%%%%%%%%%%%%%%%%%%%%%%%%%%%%%%%%%%%

\vspace{1em}
\hrule
\vspace{0.5em}

\noindent\textit{Cross-references:}
Appendix TKT (METHOD-TOOLKIT),
Appendix UNMAPPED_PBB (FORMAL-PROBABILITY),
Appendix UNMAPPED_STA (CORE-STAGE),
Appendix PAP1 (FORMAL-SEGMENT),
Appendix UNMAPPED_HOW (CORE-HOW).

\noindent\textit{Chapters:} Chapter 15 (WEC), Chapter 17 (Intervention Foundations).


% =============================================================================
% END OF APPENDIX III
% =============================================================================
